\documentclass[11pt, a4paper]{article}
\usepackage[a4paper, top=2.5cm, bottom=2.5cm, left=2cm, right=2cm]{geometry}
\usepackage[portuguese]{babel}

\usepackage{amsmath,amsthm,amssymb,amsfonts,setspace}
\usepackage{enumitem}
\usepackage{tcolorbox}

\textwidth 7.0 truein
\oddsidemargin -0.25in   %left-hand edge
\evensidemargin -0.5 truein  %right-hand edge
\topmargin -0.85in      %top of paper to top of head, pulls whole unit
\textheight 9.5in

%%%% In most cases you won't need to edit anything above this line %%%%

\begin{document}
\hfill Matheus Ferreira Souza — UC24203003

\hfill Estatística

\hfill Lista de exercício 5

\begin{tcolorbox}[colback=gray!10, colframe=black, title=Observação:]
Alguns valores críticos da distribuição Normal Padrão, como $z=1,96$ para
95\% de confiança, foram apresentados diretamente por serem valores bastante
conhecidos e recorrentes em exercícios da disciplina. Dessa forma, evitam-se
consultas repetitivas à tabela e cálculos desnecessários ao longo da resolução.
\end{tcolorbox}

\vspace{1em}

\section{Exercício 1}
\begin{enumerate}[label=\alph*)]
    \item \textbf{Resposta:}\\
    O parâmetro de interesse é a média populacional do tempo de execução do algoritmo.
    
    \item \textbf{Resposta:}\\
    O estimador utilizado é a média amostral $(\bar{X})$.
    
    \item \textbf{Resposta:}\\
    A estimativa média populacional é a própria média da amostra: $\bar{x} = 2,8s$.
\end{enumerate}                   

\section{Exercício 2}
\begin{enumerate}[label=\alph*)]
    \item \textbf{Resposta:}\\
    Ambos os estimadores possuem a propriedade de imparcialidade:\\
    $E(\hat{\theta}_1) = \mu$\\$E(\hat{\theta_2}) = \mu$\\
    \\Ou seja, os dois estimadores produzem o mesmo valor correto.
    
    \item \textbf{Resposta:}\\
    O $\hat{\theta_2}$ deve ser preferido.
    
    \item \textbf{Resposta:}\\
    Embora ambos sejam imparciais, o segundo estimador possui menor variância:\\
    $Var(\hat{\theta}_1) = 25$\\x\\$Var(\hat{\theta}_2) = 9$\\
    \\Menor variância == valor verdadeiro mais próximo de $\mu$.
    
\end{enumerate}

\section{Exercício 3}
\textbf{Resposta:}
\begin{align*}
    &4,2 \pm 1,96 * (1,5/\sqrt{100})\\
    &LS = 4,2 + 0,294 \approx 4,49\\
    &LI = 4,2 - 0,294 \approx 3,91\\
    \\&IC[3,91;4,49]\\
\end{align*}

Com 95\% de confiança, a média populacional do tempo diário de uso da plataforma está entre 3,91 e 4,49 horas.

\section{Exercício 4}
\begin{enumerate}[label=\alph*)]
    \item \textbf{Resposta:}\\
    Sim, pois $99\%$ está contido dentro do intervalo $IC[98,2\%\:;\:99,4\%]$
    
    \item \textbf{Resposta:}\\
    Não, o valor $97,5\%$ não é compatível com os dados porque é menor que o limite inferior do intervalo $(98,2\%)$, logo está fora do intervalo de confiança.
    
    \item \textbf{Resposta:}\\
    Um intervalo de $95\%$ de confiança significa que, se você repetir o mesmo estudo e/ou coleta de dados várias vezes sob as mesmas condições, cerca de $95\%$ desses intervalos terão o valor real e verdadeiro da população. 
\end{enumerate}

\section{Exercício 5}
\begin{enumerate}[label=\alph*)]
    \item \textbf{Resposta:}\\
    $H_0:\mu \leq 500hrs$ 
    
    \item \textbf{Resposta:}\\
    $H_1: \mu > 500hrs$
    
    \item \textbf{Resposta:}\\
    Unilateral à direita, porque o interesse é saber se o resultado é superior a 500 horas.
\end{enumerate}

\section{Exercício 6}
\begin{enumerate}[label=\alph*)]
    \item \textbf{Resposta:}\\
    $H_0: \mu = 12$
    
    \item \textbf{Resposta:}\\
    $H_1: \mu \neq 12$ 
    
    \item \textbf{Resposta:}\\
    Bilateral porque uma diferença para mais ou para menos é igualmente relevante.
    
    \item \textbf{Resposta:}\\
    A região crítica é composta por qualquer valor de $Z$ calculado que seja menor que $-1,96$ ou maior que $1,96$.
    \[
    RC = \{z < -1,96 \text{ ou } z > 1,96\}
    \]
\end{enumerate}

\section{Exercício 7}
\begin{enumerate}
    \item \textbf{Resposta:}\\
    $H_0: \mu = 2s$\\$H_1: \mu \neq 2s$
    
    \item \textbf{Resposta:}
    \begin{align*}
        &E = \frac{\sigma}{\sqrt{n}}\\
        &E = \frac{0,8}{\sqrt{64}} = \frac{0,8}{8} = 0,1s
    \end{align*}

    \begin{align*}
        &Z = \frac{\bar{x} - \mu}{E}\\
        &Z = \frac{2,3 - 2}{0,1} = \frac{0,3}{0,1} = 3
    \end{align*}
    
    \item \textbf{Resposta:}\\
    Região de rejeição: $z < -1,96$ ou $z > 1,96$\\
    Como $Z = 3 > 1,96$, rejeita-se a $H_0$.
    
    \item \textbf{Resposta:}\\
    Rejeita-se a hipótese nula ($H_0$) ao nível de significância de 5\%. Existem evidências estatísticas suficentes para afirmar que o tempo médio de carregamento é diferente de 2 segundos.
\end{enumerate}

\section{Exercício 8}
\begin{enumerate}
    \item \textbf{Resposta:}\\
    $H_0: \mu \geq 15s$\\
    $H_1: \mu < 15s$
    
    \item \textbf{Resposta:}
    \begin{align*}
        &E = \frac{\sigma}{\sqrt{n}}\\
        &E = \frac{3,6}{\sqrt{81}} = \frac{3,6}{9} = 0,4s
    \end{align*}

    \begin{align*}
        &Z = \frac{\bar{x} - \mu}{E}\\
        &Z = \frac{14,2 -  15}{0,4} = \frac{-0,8}{0,4} = -2\\[1em]
    \end{align*}
    
    \item \textbf{Resposta:}\\
    Se $\alpha = 1\%$ então $z_c = -2,33$\\
    Logo: $RC = \{z < -2,33\}$
    
    \item \textbf{Resposta:}\\
    Como $-2 > -2,33$, a estatística de teste não pertence a região crítica. Portanto, não se rejeita $H_0$.
    
    \item \textbf{Resposta:}\\
    A nível de significância de 1\% não existem evidências estatísticas suficientes para afirmar que a nova versão reduziu o tempo médio de processamento do sistema.
\end{enumerate}

\section{Exercício 9}
$t_{\alpha/2;24} = t_{0,025;24} = 2,064$

\begin{align*}
    &18,4 \pm 2,064 * (\frac{3,5}{\sqrt{25}})\\
    &LS = 18,4 + 1,4448 \approx 19,85s\\
    &LI = 18,4 - 1,4448 \approx 16,96s
\end{align*}

\[IC[16,96s\:;\:19,85s]\]\\[1em]

Com um nível de confiança de 95\%, a média populacional do tempo de processamento de dados está entre $16,96s$ e $19,85s$.

\end{document}